\thispagestyle{josfirst}

\noindent{\jostitlefont\fontsize{14pt}{18pt}\selectfont
事件驱动智能体操作系统跨层能力治理方法\par}

\vspace{7pt}
\noindent{\josauthorfont\fontsize{12pt}{16pt}\selectfont
\mbox{徐刚$^{1,3}$}，\mbox{冯骐$^{2,*}$}，\mbox{姚俊杰$^{3}$}，\mbox{肖宇$^{3}$}，\mbox{陈铭松$^{3}$}，\mbox{王江涛$^{1,3}$}，\mbox{杨世光$^{3}$}，\mbox{饶振宇$^{3}$}，\mbox{王语欣$^{3}$}\par}

\vspace{4pt}
\begingroup\fontsize{8pt}{12pt}\selectfont
\noindent$^{1}$（华东师范大学 国家可信嵌入式软件工程技术研究中心，上海）\par
\noindent$^{2}$（华东师范大学 信息化治理办公室，上海）\par
\noindent$^{3}$（华东师范大学 软件工程学院，上海）\par
\noindent 通讯作者：冯骐，E-mail：\texttt{qfeng@admin.ecnu.edu.cn}\par
\endgroup

\begin{joszhabstract}
事件驱动智能体操作系统将外部事件转化为自主任务并发起工具调用，但事件来源、任务创建与工具执行分属不同信任域，仅凭协议格式校验或工具允许集合检查，不能建立一次工具调用的完整链路授权依据。本文提出事件到智能体治理（Event-to-Agent Governance，E2AG）方法，将事件接入、任务创建和工具调用三个阶段的授权约束统一纳入任务作用域的能力生命周期。E2AG 通过来源--类型契约和目标策略完成任务准入，为获准任务签发对象绑定的短期能力，并在实际模型上下文协议（Model Context Protocol，MCP）调用抵达上游前实施第二次强制验证；统一执行证据支持授权关系复核和治理失败定位。本文给出跨层授权状态模型、完整中介安全性质及其证明概要，并在事件驱动智能体操作系统原型上实现该方法。评估包括经盲表复核的治理矩阵、消融与端到端执行、外部机制共同场景对照、跨项目历史事件回放、异构配置验证及证据故障注入。评估结果表明，在冻结治理矩阵中，完整方法放行全部预期可直接准入的事件，并对其余事件作出预期的拒绝或审批判定。消融显示移除准入或调用检查会使相应未授权调用到达上游；跨项目回放中，E2AG、完整上下文策略引擎和共享准入宿主的运行时规则均完成正常业务并阻断指定越权请求。实验验证了所测执行路径中跨层授权约束的有效性，以及该治理结构与现有策略机制的可组合性。
\end{joszhabstract}
\joskeywords{人工智能操作系统；智能体操作系统；能力治理；完整中介；执行溯源}
\josclassification{TP311}

\vspace{8pt}
\noindent{\bfseries\fontsize{10.5pt}{13pt}\selectfont
A Cross-Layer Capability Governance Method for Event-Driven Agent Operating Systems\par}
\vspace{4pt}
\noindent{\fontsize{9pt}{12pt}\selectfont
\mbox{XU Gang$^{1,3}$}, \mbox{FENG Qi$^{2,*}$}, \mbox{YAO Junjie$^{3}$}, \mbox{XIAO Yu$^{3}$}, \mbox{CHEN Mingsong$^{3}$}, \mbox{WANG Jiangtao$^{1,3}$}, \mbox{YANG Shiguang$^{3}$}, \mbox{RAO Zhenyu$^{3}$}, \mbox{WANG Yuxin$^{3}$}\par}
\vspace{3pt}
\begingroup\fontsize{7.5pt}{10pt}\selectfont
\noindent$^{1}$ (National Trusted Embedded Software Engineering Technology Research Center, East China Normal University, Shanghai, China)\par
\noindent$^{2}$ (Information Governance Office, East China Normal University, Shanghai, China)\par
\noindent$^{3}$ (School of Software Engineering, East China Normal University, Shanghai, China)\par
\endgroup

\begin{josenabstract}
Event-driven agent operating systems transform external events into tasks and tool invocations across distinct trust domains. Protocol validation or tool-allowlist checking alone does not establish authorization across this chain. This paper presents Event-to-Agent Governance (E2AG), which organizes event-admission, task-creation, and tool-invocation constraints into a task-scoped capability lifecycle. E2AG admits tasks through source--type contracts and target policies, issues object-bound short-lived capabilities, and checks actual Model Context Protocol (MCP) calls before forwarding. Unified evidence supports authorization review and failure localization. We define an authorization transition model and a complete-mediation property with a proof sketch, and implement E2AG in an agent operating system prototype. Evaluation combines a blindly reviewed governance matrix, ablations, end-to-end executions, external-mechanism comparisons, cross-project historical-event replay, heterogeneous configurations, and fault injection. On the frozen matrix, E2AG admits every directly admissible event and produces the expected denial or approval decision for the remaining events. Removing either enforcement point allows corresponding unauthorized calls to reach the upstream tool. In cross-project replay, E2AG, a complete-context policy engine, and runtime rules with shared admission all complete normal workflows and block the specified unauthorized requests. These results support the effectiveness of cross-layer authorization on the tested paths and its composability with existing policy mechanisms.
\end{josenabstract}
\josenkeywords{artificial intelligence operating system; agent operating system; capability governance; complete mediation; execution provenance}

\vspace{8pt}
